\section{Reproduction: what the original design can and cannot express}
\label{sec:repro}

Before we measure any disagreement with the original study, we first check
that we agree with its own published numbers.

\subsection{The published advantage survives cache-aware accounting}

We recompute the original's headline table under our own cost model: we take
their reported token counts and combine them with our measured per-token
tariffs (prices) and our measured cache-hit rates for the two archetypal
runtimes. The result narrows their advantage but does not reverse it\live:

\begin{center}
\begin{tabular}{@{}rrr@{}}
\toprule
Horizon $T$ & token ratio & cost ratio under caching \\
\midrule
10  & 1.6$\times$  & \textbf{0.97$\times$} \\
25  & 2.9$\times$  & 1.75$\times$ \\
50  & 5.7$\times$  & 3.42$\times$ \\
100 & 19.0$\times$ & 11.45$\times$ \\
200 & 21.3$\times$ & 12.82$\times$ \\
\bottomrule
\end{tabular}
\end{center}

We draw two conclusions from this table. First, caching cuts the advantage
roughly in half but does not reverse it: inside the original's own design the
method remains cheaper, and we state that plainly. Second, at $T{=}10$ the
ratio has already crossed 1.0, so cost parity disappears at the shortest
horizon in their own table, and the original paper does not remark on it
because its metric cannot detect it. This falsifies hypothesis A2 using only
published data, with no new runs of our own.

This estimate applies our tariffs to their token counts, and we label it as
such: the order of magnitude holds up under their reported numbers, but the
exact multipliers are not a measurement of what they actually spent.

\subsection{The property that hides re-acquisition}

In the original generator, horizon $T$ is a count of environment events, and
every runtime answers each event with exactly one model call\origin. The
number of model calls is therefore fixed by construction and is the same in
every arm (every condition being compared).

This single property decides what the evaluation can observe. The main cost
of discarding history is having to re-acquire what was discarded, and that
cost is paid in \emph{additional steps the agent takes on its own
initiative}. In a design where the step count is fixed as an input rather
than produced as an outcome, this cost has nowhere to show up: an agent that
has forgotten a fact and must re-derive it still gets exactly $T$ calls, so
the forgetting stays invisible in the record the original paper keeps.

In an agentic setting the step count is an outcome we measure, not an input
we fix. On our corpus, the state-only arm (the condition that keeps only the
summarized state and drops the transcript) takes \TurnRatioTailZero$\times$
as many turns as the control (the condition that keeps the full transcript)
on tasks \emph{both} arms solved\live. We measure this ratio only on tasks
where the method succeeded, so it is a lower bound on the overhead, not an
average that includes its failures.

\subsection{What this confirms and what it does not}

We reproduce the original's mechanics, its token claim, and the sign of its
cost advantage under caching, all \emph{inside its own design}. What we do
not confirm is that this design's conclusion carries over to agentic
execution. Both evaluations are legitimate, and they answer different
questions; only one of those questions is the one a practitioner needs
answered before putting a memory architecture into a coding agent.

We state one fidelity caveat here rather than leave it buried: our state
block is thicker per step than the figure the original reports\live, so our
absolute cost numbers are upper bounds. A thinner block would strengthen the
prompt-size reduction and weaken the cost inversion; both changes would work
against our headline result, which is why we report both.
\note{the paper-fidelity arm that isolates this is built and pre-registered;
if it runs before submission its numbers replace this paragraph}
